\documentclass[11pt]{article}
\usepackage[margin=1in]{geometry}
\usepackage{times}
\usepackage{graphicx}
\usepackage{booktabs}
\usepackage{longtable}
\usepackage{array}
\usepackage{hyperref}
\usepackage{enumitem}
\usepackage{titlesec}
\usepackage{xcolor}

\hypersetup{colorlinks=true, linkcolor=black, urlcolor=blue}
\titleformat{\section}{\large\bfseries}{\thesection}{1em}{}
\titleformat{\subsection}{\normalsize\bfseries}{\thesubsection}{1em}{}

\title{\textbf{Manufacturing Quality Intelligence System (MQIS)\\
A Real-Time Defect Detection and RAG-Based Quality Reporting Platform}}

\author{
Bharat Bhardwaj\thanks{Roll No: 1024240013, Email: bbhardwaj\_be24@thapar.edu } \and
Rishit\thanks{Roll No: 1024240030, Email: rgoel1\_be24@thapar.edu} \and
Asmita\thanks{Roll No: 1024240014, Email: aasmita\_be24@thapar.edu}
}
\date{August 16, 2026}

\begin{document}
\maketitle

\noindent\textbf{Project Proposal for UCS503}\\
\textbf{Submitted to:} Asif Jeelani\\
\textbf{Thapar Institute of Engineering and Technology}

\vspace{0.3cm}

\section{Higher-order Goal}
This project aims to improve manufacturing quality assurance by reducing the time and
cognitive effort required to detect, log, and diagnose defects on a production line. While
automated visual inspection already exists in industry, the immediate engineering objective
of this project is to build a deployable, end-to-end software pipeline that not only flags
defective parts but also \emph{explains} why they are occurring, using live analytics
combined with retrieval-augmented generation (RAG).

\section{Time-to-value}
\begin{itemize}[leftmargin=*]
    \item We will deliver meaningful increments quickly by prototyping the core detection-to-logging
    pipeline first, and validating it against a public defect dataset before adding intelligence layers.
    \item We will prioritize fast feedback over perfection by using weekly milestones, testing each
    module (detection, database, analytics, RAG) independently before integration.
    \item Success is defined by what can be reliably detected, logged, and reported in the first
    working prototype, then iteratively improved with better classification and reporting quality.
\end{itemize}

\section{Problem Statement}
Manufacturing lines in small and mid-sized facilities frequently rely on manual visual inspection
to catch defective parts. This approach faces several concrete issues:

\begin{itemize}[leftmargin=*]
    \item \textbf{Inconsistency:} Human inspectors tire and miss defects, especially over long shifts.
    \item \textbf{No root-cause visibility:} A rejected part is simply discarded; nothing connects
    the rejection to \emph{which} line, \emph{which} time window, or \emph{which} recurring failure
    pattern is responsible.
    \item \textbf{Delayed reporting:} Shift reports summarizing what failed and why are compiled
    manually, well after the fact, delaying corrective action.
    \item \textbf{Siloed knowledge:} Experienced technicians know that certain defect signatures
    map to specific causes (e.g., a scratch pattern indicating a worn clamp), but this knowledge is
    not captured anywhere the system can use it automatically.
\end{itemize}

\noindent\textbf{Why this matters now:} As production lines scale in number and throughput, even a
small delay in detecting a systemic defect (e.g., a mis-calibrated press) can result in a large
batch of rejected output before a human notices the pattern.

\section{Proposed Solution --- Software Proposition}

\subsection{Overview}
We propose \textbf{MQIS}, a web-based defect detection and quality-intelligence dashboard with
four integrated components:
\begin{itemize}[leftmargin=*]
    \item \textbf{Vision-based inspection module:} captures/loads a part image and compares it
    against a known-good reference using image differencing and contour analysis to localize and
    classify the defect (type, severity, confidence).
    \item \textbf{Event logging layer:} every inspection --- pass or reject --- is persisted with a
    timestamp, production line ID, defect type, severity, and confidence score.
    \item \textbf{Live analytics engine:} aggregates the log in real time into per-line rejection
    rates, line status flags (Normal / Warning / Critical), defect-type breakdowns, and hourly
    trend data (including the peak-failure hour).
    \item \textbf{RAG-based reporting assistant:} embeds a small factory knowledge base (defect
    causes, maintenance notes) into a vector store, retrieves the most relevant knowledge for the
    current data snapshot, and uses a local LLM to generate a grounded, human-readable shift report
    or answer free-text operator questions.
\end{itemize}

\subsection{Core Workflow}
\begin{enumerate}[leftmargin=*]
    \item An inspector (or the simulated feed) triggers an inspection for a selected production line.
    \item The vision module compares the part against a reference image, localizes any defect region,
    and classifies it by type and severity.
    \item The result is logged to the database and reflected immediately on the live dashboard.
    \item The analytics engine recomputes rejection rates, line status, and defect distributions.
    \item On request, the RAG assistant retrieves relevant factory knowledge and produces a shift
    report or answers a specific question (e.g., ``why is Line 1 rejecting more today?'').
\end{enumerate}

\subsection{Operational Constraints}
\begin{itemize}[leftmargin=*]
    \item The dashboard must remain responsive on standard laptop-class hardware (no dedicated GPU
    required for inspection; the LLM component runs locally via a lightweight quantized model).
    \item Avoid dependence on proprietary cloud vision APIs; use open, inspectable computer-vision
    techniques so the system's decisions remain explainable.
    \item Design the schema and analytics layer to be dataset-agnostic, so the same pipeline can
    later be pointed at a different defect category or product line with minimal changes.
\end{itemize}

\section{Solution Approach --- Engineering Focus}
This is an engineering systems project; the focus is on building a correct, integrated pipeline
rather than on inventing new defect-detection algorithms.

\subsection{Detection Module (Non-research)}
\begin{itemize}[leftmargin=*]
    \item Reference-based differencing: grayscale conversion, Gaussian blur, absolute difference
    against a known-good image.
    \item Thresholding and morphological operations (close/open) to isolate the defect region while
    suppressing noise.
    \item Contour extraction to obtain a bounding box for visualization, with a minimum-area filter
    to reject spurious detections.
    \item Defect classification is data-driven (based on labelled categories in the dataset) rather
    than a novel model, keeping scope realistic for the semester.
\end{itemize}

\subsection{System Architecture (3-tier)}
\begin{itemize}[leftmargin=*]
    \item \textbf{Frontend:} Streamlit-based dashboard --- inspection view, live statistics, defect
    log table, and an AI assistant panel.
    \item \textbf{Backend / logic:} Python modules for detection, analytics aggregation, and RAG
    orchestration.
    \item \textbf{Data layer:} SQLite for structured event logging; a local vector store (Chroma)
    for the factory knowledge base used in retrieval.
\end{itemize}

\subsection{CI/CD and Iteration Discipline}
\begin{itemize}[leftmargin=*]
    \item Each module (detector, database, analytics, RAG) is developed and unit-tested independently
    before integration into \texttt{app.py}.
    \item A shared Git repository with a linear commit history per feature branch, merged weekly.
    \item Lightweight automated checks (import sanity, schema validation) run before each merge to
    avoid breaking the dashboard.
\end{itemize}

\section{Evaluation Criteria}

\subsection{Primary Evaluation Metric}
\textbf{Mean Time-to-Log (MTTL):} the time between an inspection event occurring and it being
reflected in the live dashboard and database.
\begin{itemize}[leftmargin=*]
    \item Target: MTTL $\leq$ 2 seconds per inspection under normal load.
    \item Attribution: measured via internal timestamps between capture and successful database write.
\end{itemize}

\subsection{Secondary Evaluation Metrics}
\begin{itemize}[leftmargin=*]
    \item \textbf{Detection consistency:} percentage of known defective samples correctly flagged
    as rejected across the test dataset.
    \item \textbf{Report generation latency:} time taken by the RAG pipeline to produce a full shift
    report (target: under 30 seconds on local hardware).
    \item \textbf{Analytics correctness:} rejection-rate and status-flag calculations validated
    against manually computed values on a fixed sample log.
    \item \textbf{Usability:} informal feedback (1--5 scale) from 3--5 peer testers on dashboard
    clarity and report usefulness.
\end{itemize}

\subsection{Validation Plan}
Run a fixed batch of $\geq$50 inspections across three simulated production lines, compare
system-reported rejection rates and status flags against a manually tallied ground truth, and
collect qualitative feedback on the generated shift reports from 3--5 peer reviewers.

\section{Scalability}
\begin{enumerate}[leftmargin=*]
    \item \textbf{Theoretical:} The architecture decouples the detection, logging, analytics, and
    reporting layers, allowing each to scale independently (e.g., swapping SQLite for a managed
    relational database under higher throughput) without redesigning the pipeline.
    \item \textbf{Practical:} The knowledge base and vector store can be extended with additional
    factory documentation without code changes, and the detection module can be pointed at a new
    defect category by supplying a new reference image set and label mapping.
\end{enumerate}

\section{Engine Availability Heuristic}
A mature set of tools is available to implement this without building infrastructure from scratch:
\begin{itemize}[leftmargin=*]
    \item \textbf{OpenCV} --- image differencing, thresholding, contour detection.
    \item \textbf{Streamlit} --- rapid dashboard/UI development.
    \item \textbf{Pandas} --- analytics aggregation (rejection rates, groupby summaries).
    \item \textbf{SQLite} --- lightweight structured event storage.
    \item \textbf{ChromaDB + Sentence-Transformers} --- embedding and retrieval for the knowledge base.
    \item \textbf{Ollama (Llama 3.2)} --- local LLM inference for report generation, avoiding
    dependence on external API costs or connectivity.
\end{itemize}
We will use these mature, well-documented components to focus effort on system integration,
correctness, and evaluation rather than reinventing infrastructure.

\section{Project Scope and Deliverables}

\subsection{Initial Deliverable (Milestone 1)}
\begin{itemize}[leftmargin=*]
    \item Working detection module against a labelled reference dataset.
    \item SQLite logging of every inspection event.
    \item Minimal Streamlit dashboard showing the latest inspection result.
\end{itemize}

\subsection{Subsequent Deliverables}
\begin{itemize}[leftmargin=*]
    \item Full live-analytics panel (per-line status, defect-type/hourly breakdowns).
    \item RAG knowledge base integration with shift-report generation.
    \item Free-text question-answering assistant grounded in live data.
    \item Final testing, documentation, and demo packaging.
\end{itemize}

\section{Team Roles and Resources}
\begin{itemize}[leftmargin=*]
    \item \textbf{Bharat Bhardwaj} --- Vision/detection module, system integration across all
    three components, dataset preparation, and evaluation of detection consistency.
    \item \textbf{Rishit} --- Database schema, analytics engine, dashboard integration
    (Streamlit frontend).
    \item \textbf{Asmita} --- RAG pipeline (knowledge base, embeddings, LLM integration),
    documentation, and final report writing.
\end{itemize}
\noindent\textbf{Material resources:} A public defect-detection image dataset (e.g., MVTec-style
categories), standard laptop hardware, and open-source software only. No external funding is
required; all tools used are free and open-source.

\section{15-Week Timeline}

\begin{longtable}{@{}p{1.6cm}p{4.3cm}p{4.3cm}p{4.3cm}@{}}
\toprule
\textbf{Week} & \textbf{Bharat (Detection \& Integration)} & \textbf{Rishit (Data/Analytics)} & \textbf{Asmita (RAG/Reporting)} \\
\midrule
\endhead
1 & Literature/tooling survey; finalize dataset & Define DB schema for event logging & Survey RAG stack (Chroma, embeddings, Ollama) \\
2 & Set up reference-image comparison pipeline & Implement \texttt{database.py} (init, log, query) & Set up local Llama 3.2 via Ollama; test basic chat call \\
3 & Implement grayscale diff + thresholding & Test DB read/write with dummy data & Draft initial factory knowledge base (.txt notes) \\
4 & Add contour extraction + bounding box & Build minimal Streamlit shell (\texttt{app.py} skeleton) & Implement embedding pipeline for knowledge base \\
5 & Implement defect classification mapping & Wire detection output into DB logging & Implement Chroma collection build/query \\
6 & \textbf{Milestone 1:} Working detection on sample dataset & \textbf{Milestone 1:} Logging pipeline integrated & \textbf{Milestone 1:} Retrieval returns relevant chunks \\
7 & Tune thresholds; reduce false positives & Build \texttt{analytics.py}: rejection rate, by-line summary & Draft prompt template for shift report generation \\
8 & Add severity/confidence scoring & Add by-type and by-hour breakdowns & Implement \texttt{generate\_report()} end-to-end \\
9 & Visual annotation (bounding box + label overlay) & Add line-status flags (Normal/Warning/Critical) & Implement \texttt{answer\_question()} for free-text queries \\
10 & Integration testing with live dashboard & Live stats panel in Streamlit (metrics, bar chart) & Integrate AI assistant panel into dashboard UI \\
11 & \textbf{Milestone 2:} Full inspection view working & \textbf{Milestone 2:} Full analytics panel working & \textbf{Milestone 2:} Full RAG panel working \\
12 & End-to-end integration (all modules combined) & End-to-end integration (all modules combined) & End-to-end integration (all modules combined) \\
13 & System testing: run $\geq$50 inspections, log accuracy & Validate analytics against manual ground truth & Evaluate report quality/latency against target \\
14 & Bug fixing; buffer week for unresolved issues & Bug fixing; buffer week for unresolved issues & Bug fixing; buffer week for unresolved issues \\
15 & Final demo prep, documentation, presentation & Final demo prep, documentation, presentation & Final demo prep, documentation, presentation \\
\bottomrule
\end{longtable}

\section{Risks and Mitigations}
\begin{itemize}[leftmargin=*]
    \item \textbf{Detection accuracy on edge cases:} lighting/angle variation may cause false
    positives; mitigate with a fixed capture setup and threshold tuning during Weeks 6--8.
    \item \textbf{LLM latency/resource constraints:} local LLM inference may be slow on limited
    hardware; mitigate by using a small quantized model and caching repeated report queries.
    \item \textbf{Integration delays:} three independently developed modules may have interface
    mismatches; mitigate with an agreed data-contract (schema/field names) fixed by Week 3, and a
    dedicated integration buffer (Weeks 12, 14).
    \item \textbf{Knowledge base quality:} a thin knowledge base leads to shallow RAG answers;
    mitigate by iteratively expanding factory notes based on gaps found during testing.
\end{itemize}

\section{Summary}
This project proposes MQIS, a deployable software platform that unifies real-time visual defect
detection, structured event logging, live analytics, and RAG-based natural-language reporting for
a manufacturing quality-control setting. It meets the course criteria by emphasizing incremental,
testable delivery across three clearly divided engineering roles, using measurable evaluation
criteria (particularly mean time-to-log and detection consistency), and remaining focused on
system integration and correctness rather than algorithmic novelty.

\end{document}

